\section{The Poisson distribution}

We say that $X$ has the Poisson distribution with
parameter $\lambda>0$, and denote
$X\sim \operatorname{Poi}(\lambda)$, if $X$ satisfies

\[
\mathrm{P}(X=k) = e^{-\lambda} \frac{\lambda^k}{k!} \qquad (k\ge 0).
\]

\noindent $\textbf{Properties:}$

\begin{enumerate}[label=\arabic*.]
\item $\operatorname{Poi}(\lambda) \mathop{\boxplus}
\operatorname{Poi}(\mu) =
\operatorname{Poi}(\lambda+\mu)$.

\item The Poisson distribution is the limit of the binomial
distributions $\operatorname{Bin}(n,p)$ where the
number $n$ of experiments tends to infinity and the
probability $p$ of success in each individual
experiment goes to $0$ in such a way that the mean
number $np$ of successes stays fixed. More precisely,
if $X\sim \operatorname{Poi}(\lambda)$ and for each
$n$, $W_n$ is a r.v. with distribution
$\operatorname{Bin}(n,\lambda/n)$, then

\[
\mathrm{P}(W_n = k) \xrightarrow[n\to\infty]{} \mathrm{P}(X=k) \qquad (k\ge 0).
\]

\noindent (This is known as the $\textbf{law of rare events}$;
see Section \ref{sec:poisson-limit} for the proof of a
similar result that holds in much greater generality.)

\item If $X \sim \operatorname{Poi}(\lambda)$ then
$\mathbf{E} X = \lambda$, $\mathbf{V}(X)=\lambda$.
\end{enumerate}
